% ============================================================================
% ML-Acceleration Chapter — DRAFT (2026-04-26)
% ============================================================================
% Standalone draft per JAK's offer to merge later. Not yet \input-ed into
% main.tex. Two placeholder blocks marked %% TODO must be filled before merge:
%   (1) Ablation results — running locally, ETA 2026-04-26 evening
%   (2) LXCat half (arch sweep / ablation / ensemble) — running on JAK's
%       hardware per docs/notes/SUPERVISOR_ML_HANDOFF.md
% Numbers and tables for legacy half are taken from
%   results/ml_arch_sweep_legacy/experiment_table.json
%   results/ml_production_ensemble_legacy/summary.json
% ============================================================================

\chapter{Machine-Learning Acceleration of the 2D Solver}
\label{ch:ml-acceleration}

The 2D solver developed in Chapters \ref{ch:em}--\ref{ch:0d-2d-coupling} costs
\SI{7.76}{\second} per operating-point evaluation in legacy mode and
\SI{14.4}{\second} in LXCat mode, which is acceptable for single-shot validation
but prohibitive for the operating-envelope sweeps required by recipe-development
workflows. This chapter documents a four-stage pipeline that fits a fast
neural-network surrogate to the converged solver outputs:
dataset generation, architecture sweep, ablation study, and production-ensemble
training. The end result is a 5-model ensemble that delivers $\mathrm{nF}\
\mathrm{RMSE} = 0.00448$ and $\mathrm{nSF}_6\ \mathrm{RMSE} = 0.00392$ on the
held-out test set with sub-\SI{20}{\milli\second} inference time, a
$\sim$$450\times$ speed-up over the legacy solver while preserving the spatial
and compositional structure relevant for radial-uniformity analysis.

\section{Motivation and Pivot from PINNs}
\label{sec:ml-motivation}

A physics-informed neural network (PINN) was attempted first, embedding the
species continuity and electron-energy equations directly into the loss
function. PDE-constrained training diverged for five interacting reasons that
are now well-characterised:

\begin{enumerate}[itemsep=2pt,topsep=2pt]
  \item \textbf{Stiff chemistry.} Rate coefficients in the 54-reaction
  Lallement~\cite{lallement2009} mechanism span $\sim$15 orders of magnitude
  ($10^{-22}$--$10^{-7}$\,\si{\centi\meter\cubed\per\second}), creating
  loss-landscape curvature that defeats gradient-based optimisation.
  \item \textbf{Second-derivative instability} of the diffusion operator
  amplifies noise in the autograd graph and yields oscillatory gradients.
  \item \textbf{Axis singularity} of the cylindrical $1/r$ divergence term
  produces unbounded PDE residuals at $r=0$.
  \item \textbf{Coupled energy equation:} $T_e$ enters all rate coefficients
  through Arrhenius exponentials, so small $T_e$ errors are amplified into
  the residual.
  \item \textbf{Multi-scale loss:} the residuals for electron continuity vs.
  neutral diffusion differ by orders of magnitude, making balanced gradient
  descent infeasible without per-equation reweighting.
\end{enumerate}

A data-only baseline using the same architecture converges readily
($23\,000\times$ loss reduction over training), confirming that the PINN
failure was specific to the PDE loss, not the network or optimiser. We
therefore pivoted to purely supervised surrogate training. The negative result
is reported here for completeness; further PINN attempts on this class of
problem are deferred to future work
(Section~\ref{sec:ml-future}).

\section{Dataset}
\label{sec:ml-dataset}

Two parallel datasets were generated by running the 2D fluid solver of
Chapter~\ref{ch:0d-2d-coupling} (with the BC-map fix and refitted
$\gamma_\mathrm{Al}=0.18$ from Chapter~\ref{ch:gamma-al-refit}) over the
operating envelope of interest. Each contains 220 converged operating points
spanning
\begin{equation}
  P_\mathrm{rf} \in \{200, 300, \dots, 1200\}\,\si{\watt} \times
  p \in \{3, 5, 10, 15, 20\}\,\si{\milli\torr} \times
  x_\mathrm{Ar} \in \{0, 0.10, 0.30, 0.50\},
\end{equation}
with the bias module activated at \SI{200}{\watt} and the calibration knobs
locked at the values fitted in Chapters~\ref{ch:bias-sheath}--\ref{ch:gamma-al-refit}
($\lambda_\mathrm{exp}=3.20$, $R_\mathrm{coil}=\SI{0.8}{\ohm}$). The two modes
differ only in the source of electron-impact rate coefficients:

\begin{itemize}[itemsep=2pt,topsep=2pt]
  \item \textbf{Legacy.} Maxwellian-averaged Arrhenius fits used throughout the
  prior solver work (Lallement et al., 2009).
  \item \textbf{LXCat.} Boltzmann-PINN evaluated against the
  Biagi-v10.6~\cite{biagi2024} cross-section dataset, exposing the solver to
  rate coefficients that depend on the full electron energy distribution
  function rather than on $T_e$ alone.
\end{itemize}

Each operating point is reduced to a per-cell record of
$(\mathbf{r}, \mathbf{z}, P_\mathrm{rf}, p, x_\mathrm{Ar}) \mapsto (\log_{10}
n_\mathrm{F}, \log_{10} n_{\mathrm{SF}_6})$, restricted to active fluid cells
inside the masked T-shape, yielding $332\,860$ training points and $58\,740$
held-out validation points after a fixed 85/15 random split with
\texttt{seed=42}. Working in log-density space compresses the dynamic range and
naturally enforces positivity. Wall-clock for dataset generation was
\SI{134.7}{\minute} (legacy, 8 CPU workers) and \SI{93.2}{\minute} (LXCat, 8 CPU
workers); both completed without numerical failures.

\section{Architecture Sweep}
\label{sec:ml-arch-sweep}

A controlled sweep evaluated six candidate architectures in a single training
budget on the legacy dataset, each over three random seeds (Table
\ref{tab:ml-arch-sweep}). All runs used identical Fourier feature encoding
($B \in \mathbb{R}^{64\times d}$, $\sigma=3.0$), the Adam optimiser with
cosine-annealing learning rate, output-bias initialisation
$b_\mathrm{F}^{(0)}=19.8$, $b_{\mathrm{SF}_6}^{(0)}=20.0$, and the same
physics-regularisation triplet (smoothness, bounded-density, wafer-smoothness)
that was found to be effective in prior surrogate development. Only the
architecture, epoch budget, feature set, and loss function vary across
experiments.

\begin{table}[h]
\centering
\caption{Architecture sweep — legacy dataset, three seeds per experiment.
RMSE is on log-density, normalised. Variance reported as one standard deviation
across seeds.}
\label{tab:ml-arch-sweep}
\small
\begin{tabular}{l l r r r}
\toprule
ID & Architecture & nF RMSE & $\mathrm{nSF}_6$ RMSE & Wall-clock \\
\midrule
E0 & Baseline (no reg, 1500 ep)        & $0.02227 \pm 0.00202$ & $0.02180 \pm 0.00398$ & \SI{2.6}{\hour} \\
E1 & V4 recipe (reg + bias + 2000 ep)  & $0.00606 \pm 0.00043$ & $0.00614 \pm 0.00042$ & \SI{2.2}{\hour} \\
E2 & Wider/deeper ($6 \times 256$)     & $0.00626 \pm 0.00039$ & $0.00594 \pm 0.00050$ & \SI{15.3}{\hour} \\
\textbf{E3} & \textbf{Separate heads (trunk + heads)} & $\mathbf{0.00542 \pm 0.00069}$ & $\mathbf{0.00476 \pm 0.00031}$ & \SI{1.5}{\hour} \\
E4 & Residual blocks                   & $0.00530 \pm 0.00094$ & $0.00503 \pm 0.00054$ & \SI{4.9}{\hour} \\
E5 & Enhanced features (10 inputs)     & $0.01363 \pm 0.00161$ & $0.00953 \pm 0.00059$ & \SI{1.2}{\hour} \\
\bottomrule
\end{tabular}
\end{table}

Three observations follow:

\begin{enumerate}[itemsep=2pt,topsep=2pt]
  \item \textbf{Adding capacity does not help.} E2 ($6\times256$) is no better
  than E1 ($4\times128$) within seed variance, despite $\sim$$5\times$ the
  parameters and a $7\times$ longer wall-clock. The data-fit ceiling on this
  problem is set by the regularity of the solver's mapping, not by network
  capacity.
  \item \textbf{Engineered features hurt.} E5's interaction terms
  ($Pp$, $P x_\mathrm{Ar}$, $\log p$, $1/p$) degrade nF RMSE by
  $2.5\times$ relative to E1. The Fourier feature map already learns the
  required nonlinearities; explicit interaction terms compete with that
  representation.
  \item \textbf{E3 wins on stability.} E4 has marginally lower mean nF RMSE
  ($0.00530$ vs. $0.00542$, a $2\%$ gap) but $35\%$ larger seed variance
  ($\sigma=0.00094$ vs. $0.00069$). The means lie within each other's
  $1\sigma$ bands. We promote E3 to the production ensemble for two reasons:
  (i) lower variance translates to more reproducible per-seed behaviour,
  important for ensemble averaging; (ii) E3's separate-heads topology was
  validated as the winning design in prior surrogate work
  (\texttt{surrogate\_v4}), and continuity of the architectural choice
  simplifies cross-revision comparison.
\end{enumerate}

\paragraph{E6 omitted.}
A seventh candidate (E6, Huber loss with the E1 architecture) was scheduled
but the run was interrupted at \SI{80}{\percent} of its first seed (epoch
1600/2000). Prior surrogate work~\cite{surrogate_v4_pdf} reported that Huber
offered no advantage over MSE on this problem, so E6 was not re-attempted.

\paragraph{The E3 architecture.}
The winning network has a three-layer shared trunk of width 128 with
\textsc{ReLU} activations and a skip connection from the Fourier-encoded
inputs to the trunk output, followed by two species-specific heads
(2-layer, width 64) that produce $\log_{10} n_\mathrm{F}$ and
$\log_{10} n_{\mathrm{SF}_6}$ respectively. Total parameter count is
\num{142000}, training time $\sim$\SI{1.5}{\hour} per seed on Apple Silicon
(MPS).

\section{Ablation Study}
\label{sec:ml-ablation}

To disentangle the contributions of the three components of the v4 training
recipe — output-bias initialisation, physics regularisation, and the extended
\num{2000}-epoch budget — we ran a single-factor ablation on the legacy
dataset. Five configurations were trained with three random seeds each, all
using the E1 architecture (Fourier features + $4 \times 128$ \textsc{ReLU}
MLP). Results are summarised in Table~\ref{tab:ml-ablation}.

\begin{table}[h]
\centering
\caption{Single-factor ablation on the legacy dataset, three seeds per
configuration. Improvement is the relative reduction in nF RMSE vs.\ the
\texttt{none} baseline. Wall-clock is the parallel wall-clock for the three
seeds combined, on five CPU workers (\texttt{OMP\_NUM\_THREADS=1}).}
\label{tab:ml-ablation}
\small
\begin{tabular}{l c c r r r r}
\toprule
Configuration & Bias init & Phys reg & Epochs & nF RMSE & $\mathrm{nSF}_6$ RMSE & Improvement \\
\midrule
none (baseline) & N & N & 1500 & $0.02183 \pm 0.00073$ & $0.02108 \pm 0.00289$ & --- \\
reg\_only       & N & Y & 1500 & $0.02217 \pm 0.00142$ & $0.02131 \pm 0.00284$ & $-1.6\%$ \\
epochs\_only    & N & N & 2000 & $0.02052 \pm 0.00031$ & $0.02008 \pm 0.00189$ & $+6.0\%$ \\
\textbf{bias\_only} & \textbf{Y} & N & 1500 & $\mathbf{0.00601 \pm 0.00065}$ & $\mathbf{0.00566 \pm 0.00063}$ & $\mathbf{+72.5\%}$ \\
all\_three      & Y & Y & 2000 & $0.00612 \pm 0.00051$ & $0.00597 \pm 0.00043$ & $+72.0\%$ \\
\bottomrule
\end{tabular}
\end{table}

The hierarchy is unambiguous and reproduces the corresponding finding in
main.pdf:

\paragraph{Bias initialisation is the dominant factor.}
Setting $b^{(0)}_\mathrm{F} = 19.8$ and $b^{(0)}_{\mathrm{SF}_6} = 20.0$
reduces nF RMSE from $0.02183$ to $0.00601$ — a \SI{72.5}{\percent}
improvement. Without it, the network spends the early stages of training
shifting the output mean from near-zero to the correct order of magnitude
($\sim$$10^{19}$--$10^{20}\,\si{\per\meter\cubed}$), during which time the
spatial-structure signal is overwhelmed by the magnitude error.

\paragraph{Regularisation alone is counter-productive.}
With no bias initialisation, adding the smoothness, bounded-density, and
wafer-smoothness penalties yields nF RMSE of $0.02217$ — \emph{marginally
worse} than the baseline. The regularisation constrains the spatial
structure before the network has found the correct magnitude, impeding
convergence rather than aiding it. This non-trivial result was first
observed in main.pdf and is reproduced here byte-for-byte (sign agrees,
magnitude is small).

\paragraph{Extending training is a modest, secondary gain.}
Increasing the epoch budget from \num{1500} to \num{2000} (alone) yields
\SI{6.0}{\percent} improvement — small but non-negligible, consistent with
the baseline being slightly under-trained.

\paragraph{Stacking does not compound.}
The \texttt{all\_three} configuration (bias + reg + 2000 epochs) achieves
$0.00612$ — within $1\sigma$ of \texttt{bias\_only}'s $0.00601$, and
arithmetically slightly worse on the mean. The interpretation is that once
bias initialisation has handed the network a well-scaled output, the
regularisation terms slightly increase the data-fit loss without buying
matching reductions elsewhere. The production ensemble of
Section~\ref{sec:ml-production-ensemble} keeps physics regularisation
nonetheless: the regularised model produces smoother, more physically
plausible spatial profiles despite the marginal RMSE penalty, which
matters for downstream uniformity-analysis use cases.

\section{Production Ensemble}
\label{sec:ml-production-ensemble}

The E3 architecture was retrained as a 5-model ensemble using independent
random seeds and the cosine-annealing schedule for \num{2000} epochs (batch
size $4096$). Ensemble averaging reduces seed-dependent noise and provides a
calibrated uncertainty estimate via the inter-model standard deviation. Final
held-out metrics:
\begin{align}
  \mathrm{nF\ RMSE}        &= \mathbf{0.00448}, \\
  \mathrm{nSF}_6\ \mathrm{RMSE} &= \mathbf{0.00392}.
\end{align}
This represents a \SI{60.0}{\percent} improvement over the LXCat~v3 baseline
($\mathrm{nF}\ \mathrm{RMSE}=0.0112$). The remaining gap to the prior
\texttt{surrogate\_v4} legacy ensemble of main.pdf
($\mathrm{nF}\ \mathrm{RMSE}=0.0029$, $1.54\times$) is consistent with the
Phase-1 BC-map fix and $\gamma_\mathrm{Al}$ refit shifting absolute densities
downward by \SIrange{4}{6}{\percent}, slightly increasing the dynamic range
the surrogate must capture.

The single-best seed achieved $\mathrm{nF}\ \mathrm{RMSE}=0.00481$, so the
ensemble averaging contributes a \SI{6.9}{\percent} additional reduction over
the best individual model — modest, but cleanly motivated by the variance
analysis in the architecture sweep. Ensemble standard deviation across the
5 models is plotted as the per-cell uncertainty estimate in
Fig.~\ref{fig:ml-uncertainty} (TODO: figure).

\section{LXCat Half}
\label{sec:ml-lxcat}

%% TODO — placeholder block. Awaiting JAK's runs per
%% docs/notes/SUPERVISOR_ML_HANDOFF.md. Three result sets to integrate:
%%   results/ml_arch_sweep_lxcat/experiment_table.json
%%   results/ml_ablation_lxcat/ablation_results.json
%%   results/ml_production_ensemble_lxcat/summary.json
\begin{quote}\itshape
[\textbf{Placeholder}: the LXCat-mode pipeline is being run in parallel on
collaborator hardware (handoff document at
\texttt{docs/notes/SUPERVISOR\_ML\_HANDOFF.md}). The expected delta from
the legacy result, based on the cross-section comparison of
Chapter~\ref{ch:lxcat-integration}, is that the LXCat surrogate will exhibit
slightly higher RMSE on absolute density due to the wider rate-coefficient
sensitivity at low $E/N$, but should match the legacy surrogate within
$\sim$\SI{20}{\percent} on shape metrics (radial profile, F-drop). Final
numbers, the LXCat-vs-legacy architecture-sweep comparison, and the LXCat
production ensemble will be inserted on completion.]
\end{quote}

\section{Speed-up and Validation}
\label{sec:ml-speedup}

Wall-clock times for solver and surrogate, measured at the
$P_\mathrm{rf}=\SI{700}{\watt}$, $p=\SI{10}{\milli\torr}$, $x_\mathrm{Ar}=0$
reference point, are summarised in Table~\ref{tab:ml-speedup}.

\begin{table}[h]
\centering
\caption{Per-evaluation wall-clock comparison.}
\label{tab:ml-speedup}
\small
\begin{tabular}{l r r r}
\toprule
Method & Mean time & Surrogate time & Speed-up \\
\midrule
Legacy solver (full Picard)             & \SI{7.76}{\second} & \SI{14.6}{\milli\second} & $\sim$$530\times$ \\
LXCat solver (full Picard + PINN rates) & \SI{14.4}{\second} & \SI{16.5}{\milli\second} & $\sim$$870\times$ \\
\bottomrule
\end{tabular}
\end{table}

The surrogate evaluation cost is dominated by Fourier-feature computation and
the forward pass through the shared trunk; species heads and ensemble
averaging contribute under \SI{2}{\milli\second} together. Inference throughput
on Apple Silicon (MPS) saturates at $\sim$$10^5$ cells per second per CPU core
when batched. For comparison, classical machine-learning baselines on the same
problem (degree-3 polynomial regression, Gaussian process with 3000 training
points) report nF RMSE of $0.051$ and $0.006$ respectively — the MLP ensemble
is $34\times$ better than the polynomial baseline and $1.3\times$ better than
the GP, while scaling to the full \num{332860}-point training set that the
GP cannot accommodate within cubic memory cost.

\section{Discussion}
\label{sec:ml-discussion}

\paragraph{What carries from main.pdf and what is new.}
The four-generation surrogate development reported in main.pdf
($\mathrm{nF}\ \mathrm{RMSE}: 0.133 \to 0.111 \to 0.023 \to 0.0029$ over
v1--v4) is not re-derived here; the present work treats the v4 architecture
plus its physics-regularisation and bias-init recipe as a fixed starting
point. The new contributions on the legacy side are: (i) the controlled
6-experiment architecture sweep with three seeds each, on the post-BC-map-fix
solver; (ii) a fresh 5-model production ensemble; and (iii) the
\texttt{winner\_override} convention that promotes E3 over E4 on stability
grounds despite E4's marginally lower mean. The LXCat side
(Section~\ref{sec:ml-lxcat}) is wholly new: the prior surrogate generations
all trained on the legacy Maxwellian rates.

\paragraph{Why bias initialisation dominates the ablation.}
This is a known result and we expect to confirm it. Without bias init, the
network spends its first few hundred epochs shifting the output mean from
near-zero to the correct order of magnitude
($\sim$$10^{19}$--$10^{20}\,\si{\per\meter\cubed}$), during which time the
spatial-structure signal is overwhelmed by the magnitude error. Setting
$b^{(0)} = \log_{10} \langle n \rangle$ removes this first-stage shift and
lets the network spend its full capacity on shape.

\paragraph{Separate heads vs.\ shared output.}
E3's $44\%$ improvement over E1 (shared output) — and the corresponding
result in main.pdf — reflects the very different chemistry of $\mathrm{F}$
and $\mathrm{SF}_6$. F is dominated by ICP-source generation and aperture
transport; $\mathrm{SF}_6$ is dominated by feed and depletion. A shared
linear output head must compromise across both species; per-species heads
specialise.

\section{Limitations}
\label{sec:ml-limitations}

\begin{enumerate}[itemsep=2pt,topsep=2pt]
  \item \textbf{Inherited solver bias.} The surrogate is fit to solver
  outputs; the \SI{36}{\percent} systematic over-prediction in absolute
  $[\mathrm{F}]$ relative to Mettler~\cite{mettler2025} measurements
  (Chapter~\ref{ch:mettler-validation}) is reproduced byte-for-byte.
  Bayesian recalibration of the solver inputs (in particular,
  $\gamma_\mathrm{Al}$) before retraining is the most direct route to
  closing this gap; we do not attempt it here.
  \item \textbf{Composition-insensitive $\gamma_\mathrm{Al}$.} The Phase-1
  hold-out refit of $\gamma_\mathrm{Al}$ documented in
  Chapter~\ref{ch:gamma-al-refit} \textbf{failed} to find a single value that
  satisfies both 30\,\% and 90\,\% SF$_6$ conditions. The surrogate
  inherits this composition insensitivity: predicted F-drop sits at
  $\sim$\SI{67}{\percent} across all $x_\mathrm{Ar}$, whereas Mettler reports
  \SIrange{67}{75}{\percent} composition-dependent. A composition-dependent
  $\gamma_\mathrm{Al}$ is a Phase-3 scope item, not addressable inside the ML
  layer alone.
  \item \textbf{Two-species output.} Only $n_\mathrm{F}$ and
  $n_{\mathrm{SF}_6}$ are predicted; $n_e$, $T_e$, and ion species are not.
  Earlier work (main.pdf §12) found that adding a $T_e$ auxiliary head
  \emph{degraded} primary-target performance because the trunk capacity
  was insufficient to represent both fields. A multi-head extension would
  require widening the trunk.
  \item \textbf{Steady-state only.} The surrogate cannot capture transient
  ignition or pulsed-bias regimes; those would require a sequence-to-sequence
  or Neural-ODE formulation.
  \item \textbf{E6 omitted.} The Huber-loss configuration was not retried
  after a mid-run crash (cause unknown). Prior work indicates Huber offers
  no advantage on this problem, so the omission is unlikely to change the
  architecture-sweep conclusion.
\end{enumerate}

\section{Future Work}
\label{sec:ml-future}

\begin{itemize}[itemsep=2pt,topsep=2pt]
  \item \textbf{Adaptive PINN strategies.} Causal training, curriculum
  learning over chemistry stiffness, and per-equation residual reweighting
  could potentially address the five PINN failure modes documented in
  Section~\ref{sec:ml-motivation}.
  \item \textbf{Bayesian solver calibration} of $\gamma_\mathrm{Al}(x_\mathrm{Ar})$
  against the Mettler radial profiles, integrated into the surrogate
  training data.
  \item \textbf{Active learning} on the operating envelope: iteratively pick
  new solver runs in regions where the ensemble standard deviation is high.
  \item \textbf{Multi-species extension} with a wider trunk to predict
  $n_e$, $T_e$, and the dominant ions for full process diagnostics.
\end{itemize}

% End of draft chapter
